% GENERATED FILE. Edit canonical lessons, then run npm run generate:books.
% canonical-source-hash: fnv1a64:624f6b52869736ff
% canonical-lessons: SA-C306-sibiram, SA-C306-asrayah, SA-C306-danam, SA-C306-saranarthi, SA-C306-santvayati

\chapter{The Relief Camp}
\label{ch:sa-the-relief-camp}
\hlchaptermodality{306}

\begin{chapteropening}
\textbf{By the end of this chapter:} \emph{I can describe life in a relief camp.}

Complete the last lesson of chapter 306: I can describe life in a relief camp.
\end{chapteropening}

\section[\texorpdfstring{\sk{शिबिरम्} (śibiram)}{शिबिरम् (śibiram)}]{\sk{शिबिरम्} (śibiram) — a camp; a tent}

\label{lesson:SA-C306-sibiram}



\begin{quote}
Before the new one: say the Sanskrit for brings down, then the Sanskrit for saves.
\end{quote}

\subsection*{You'll want to know: \sk{शिबिरम्}}
\textbf{\sk{शिबिरम्}} — \emph{śibiram} — \textquotedblleft{}a camp; a tent\textquotedblright{}.

\begin{quote}
\sk{आर्ताः} \sk{जनाः} \sk{शिबिरे} \sk{वसन्ति।}
\end{quote}

\begin{quote}
\emph{ārtāḥ janāḥ śibire vasanti.}
\end{quote}

\begin{quote}
People in distress live in the camp.
\end{quote}

A neuter noun.

\begin{cousinweb}
\textbf{\sk{शिबिरम्}} (\emph{śibiram}) MW gives \emph{śibira} as a royal camp, a tent, any tent, without a derivation, so its origin is uncertain (MW).
\end{cousinweb}

\begin{grammarlens}[title={What you've built}]
One more word for telling what happened.
\end{grammarlens}

\subsection*{Your turn}
\begin{itemize}
\raggedright
  \item \emph{Say it:} \emph{śibiram}
  \item \emph{Say it:} \emph{śibiram}, once more
  \item \emph{Say it:} \emph{ārtāḥ janāḥ śibire vasanti.} — People in distress live in the camp.
  \item \emph{Recall:} say the Sanskrit for brings down, then the Sanskrit for saves, then say \emph{śibiram} again
\end{itemize}

\begin{tcolorbox}[breakable,colback=teal!4,colframe=teal!35!black,title={Before you move on}]
What does \sk{शिबिरम्} mean? (\textquotedblleft{}A camp; a tent\textquotedblright{}.) Say it once more.
\end{tcolorbox}
\section[\texorpdfstring{\sk{आश्रयः} (āśrayaḥ)}{आश्रयः (āśrayaḥ)}]{\sk{आश्रयः} (āśrayaḥ) — shelter; refuge}

\label{lesson:SA-C306-asrayah}



\begin{quote}
Before the new one: say the Sanskrit for saves, then the Sanskrit for a camp.
\end{quote}

\subsection*{You'll want to know: \sk{आश्रयः}}
\textbf{\sk{आश्रयः}} — \emph{āśrayaḥ} — \textquotedblleft{}shelter; refuge\textquotedblright{}.

\begin{quote}
\sk{वृद्धः} \sk{शिबिरे} \sk{आश्रयं} \sk{याचते।}
\end{quote}

\begin{quote}
\emph{vṛddhaḥ śibire āśrayaṃ yācate.}
\end{quote}

\begin{quote}
The old man asks for shelter in the camp.
\end{quote}

A masculine noun.

\begin{cousinweb}
\textbf{\sk{आश्रयः}} (\emph{āśrayaḥ}) is built from \emph{ā} and the root √śri, \textquotedblleft{}to lean on, resort to\textquotedblright{}; MW gives \emph{āśraya} as that on which anything rests or depends, a refuge (MW).
\end{cousinweb}

\begin{grammarlens}[title={What you've built}]
One more word for telling what happened.
\end{grammarlens}

\subsection*{Your turn}
\begin{itemize}
\raggedright
  \item \emph{Say it:} \emph{āśrayaḥ}
  \item \emph{Say it:} \emph{āśrayaḥ}, once more
  \item \emph{Say it:} \emph{vṛddhaḥ śibire āśrayaṃ yācate.} — The old man asks for shelter in the camp.
  \item \emph{Recall:} say the Sanskrit for saves, then the Sanskrit for a camp, then say \emph{āśrayaḥ} again
\end{itemize}

\begin{tcolorbox}[breakable,colback=teal!4,colframe=teal!35!black,title={Before you move on}]
What does \sk{आश्रयः} mean? (\textquotedblleft{}Shelter; refuge\textquotedblright{}.) Say it once more.
\end{tcolorbox}
\section[\texorpdfstring{\sk{दानम्} (dānam)}{दानम् (dānam)}]{\sk{दानम्} (dānam) — a donation; giving}

\label{lesson:SA-C306-danam}



\begin{quote}
Before the new one: say the Sanskrit for a camp, then the Sanskrit for shelter.
\end{quote}

\subsection*{You'll want to know: \sk{दानम्}}
\textbf{\sk{दानम्}} — \emph{dānam} — \textquotedblleft{}a donation; giving\textquotedblright{}.

\begin{quote}
\sk{जनाः} \sk{शिबिराय} \sk{अन्नस्य} \sk{दानं} \sk{प्रयच्छन्ति।}
\end{quote}

\begin{quote}
\emph{janāḥ śibirāya annasya dānaṃ prayacchanti.}
\end{quote}

\begin{quote}
People hand over food to the camp as a donation.
\end{quote}

A neuter noun.

\begin{cousinweb}
\textbf{\sk{दानम्}} (\emph{dānam}) comes from the root √dā, \textquotedblleft{}to give\textquotedblright{}, as in \emph{dadāti}; MW gives \emph{dāna} as the act of giving and a gift (MW).
\end{cousinweb}

\begin{grammarlens}[title={What you've built}]
One more word for telling what happened.
\end{grammarlens}

\subsection*{Your turn}
\begin{itemize}
\raggedright
  \item \emph{Say it:} \emph{dānam}
  \item \emph{Say it:} \emph{dānam}, once more
  \item \emph{Say it:} \emph{janāḥ śibirāya annasya dānaṃ prayacchanti.} — People hand over food to the camp as a donation.
  \item \emph{Recall:} say the Sanskrit for a camp, then the Sanskrit for shelter, then say \emph{dānam} again
\end{itemize}

\begin{tcolorbox}[breakable,colback=teal!4,colframe=teal!35!black,title={Before you move on}]
What does \sk{दानम्} mean? (\textquotedblleft{}A donation; giving\textquotedblright{}.) Say it once more.
\end{tcolorbox}
\section[\texorpdfstring{\sk{शरणार्थी} (śaraṇārthī)}{शरणार्थी (śaraṇārthī)}]{\sk{शरणार्थी} (śaraṇārthī) — a refugee}

\label{lesson:SA-C306-saranarthi}



\begin{quote}
Before the new one: say the Sanskrit for shelter, then the Sanskrit for a donation.
\end{quote}

\subsection*{You'll want to know: \sk{शरणार्थी}}
\textbf{\sk{शरणार्थी}} — \emph{śaraṇārthī} — \textquotedblleft{}a refugee\textquotedblright{}.

\begin{quote}
\sk{शरणार्थी} \sk{शिबिरे} \sk{भोजनं} \sk{प्राप्नोति।}
\end{quote}

\begin{quote}
\emph{śaraṇārthī śibire bhojanaṃ prāpnoti.}
\end{quote}

\begin{quote}
The refugee gets food in the camp.
\end{quote}

A masculine noun.

\begin{cousinweb}
\textbf{\sk{शरणार्थी}} (\emph{śaraṇārthī}) is built from \emph{śaraṇa}, \textquotedblleft{}refuge\textquotedblright{}, and \emph{arthin}, \textquotedblleft{}seeking\textquotedblright{}; MW gives \emph{śaraṇārthin} as seeking refuge, and a refugee is its modern use as a noun (MW).
\end{cousinweb}

\begin{grammarlens}[title={What you've built}]
One more word for telling what happened.
\end{grammarlens}

\subsection*{Your turn}
\begin{itemize}
\raggedright
  \item \emph{Say it:} \emph{śaraṇārthī}
  \item \emph{Say it:} \emph{śaraṇārthī}, once more
  \item \emph{Say it:} \emph{śaraṇārthī śibire bhojanaṃ prāpnoti.} — The refugee gets food in the camp.
  \item \emph{Recall:} say the Sanskrit for shelter, then the Sanskrit for a donation, then say \emph{śaraṇārthī} again
\end{itemize}

\begin{tcolorbox}[breakable,colback=teal!4,colframe=teal!35!black,title={Before you move on}]
What does \sk{शरणार्थी} mean? (\textquotedblleft{}A refugee\textquotedblright{}.) Say it once more.
\end{tcolorbox}
\section[\texorpdfstring{\sk{सान्त्वयति} (sāntvayati)}{सान्त्वयति (sāntvayati)}]{\sk{सान्त्वयति} (sāntvayati) — comforts; consoles}

\label{lesson:SA-C306-santvayati}



\begin{quote}
Before the new one: say the Sanskrit for a donation, then the Sanskrit for a refugee.
\end{quote}

\subsection*{You'll want to know: \sk{सान्त्वयति}}
\textbf{\sk{सान्त्वयति}} — \emph{sāntvayati} — \textquotedblleft{}comforts; consoles\textquotedblright{}.

\begin{quote}
\sk{माता} \sk{बालिकां} \sk{सान्त्वयति।}
\end{quote}

\begin{quote}
\emph{mātā bālikāṃ sāntvayati.}
\end{quote}

\begin{quote}
The mother comforts the girl.
\end{quote}

Past (she …): \sk{सा} \sk{बालिकां} \sk{सान्त्वितवती}, \emph{sā bālikāṃ sāntvitavatī}.

\begin{cousinweb}
\textbf{\sk{सान्त्वयति}} (\emph{sāntvayati}) is built on \emph{sāntva}, \textquotedblleft{}kind words, consolation\textquotedblright{}: MW treats \emph{sāntvayati} as rather a verb made from that noun (MW).
\end{cousinweb}

\begin{grammarlens}[title={What you've built}]
One more word for telling what happened.
\end{grammarlens}

\subsection*{Your turn}
\begin{itemize}
\raggedright
  \item \emph{Say it:} \emph{sāntvayati}
  \item \emph{Say it:} \emph{sāntvayati}, once more
  \item \emph{Say it:} \emph{mātā bālikāṃ sāntvayati.} — The mother comforts the girl.
  \item \emph{Recall:} say the Sanskrit for a donation, then the Sanskrit for a refugee, then say \emph{sāntvayati} again
\end{itemize}

\begin{tcolorbox}[breakable,colback=teal!4,colframe=teal!35!black,title={Before you move on}]
What does \sk{सान्त्वयति} mean? (\textquotedblleft{}Comforts; consoles\textquotedblright{}.) Say it once more.
\end{tcolorbox}
